\section{Discharging Form of the \(g\)-Calculus}
\label{sec:strategy2-discharging-calculus}

This appendix proves the additive identities used in
Section~\ref{sec:strategy2-discharging} and records exactly which parts of
Strategy~2 become genuine discharging arguments.  It also proves the
obstructions to a fixed, state-independent diagonal charge.  No theorem below
changes the exact transfer maps or their geometric domains.

\subsection{Release, subsidy, and actual capacity slack}

Recall
\[
 F_c(a)=\widehat g_c(1-a),
 \qquad
 G_c(a)=\widehat g_c^\vee(a)=1-F_c(a).
\]

\begin{definition}[Release and subsidy]
\label{def:appendix-release-subsidy}
For \(0\le a,c\le1\), define
\begin{equation}
 \mathfrak d_c(a):=G_c(a)-a=1-a-F_c(a).
 \label{eq:appendix-release}
\end{equation}
For an empty or closed interval \(J\subseteq[0,1]\), define
\begin{equation}
 \mathfrak t_J(p):=e_J(p)-p.
 \label{eq:appendix-subsidy}
\end{equation}
\end{definition}

Since \(F_c(a)\le1-a\), one has \(\mathfrak d_c(a)\ge0\).  Since
\([0,p]\) belongs to the component used in the definition of \(e_J(p)\), one
has \(\mathfrak t_J(p)\ge0\).

\begin{proposition}[Exact local balance law]
\label{prop:appendix-local-balance}
For every \(a,c\in[0,1]\) and every allowed interval \(J\),
\begin{equation}
 \widehat g_{c,J}^\vee(a)
 =a+\mathfrak d_c(a)-\mathfrak t_J(F_c(a)).
 \label{eq:appendix-local-balance}
\end{equation}
If an actual nonsupercritical V triangle \(T\) satisfies
\(A(T)\ge a\), \(C(T)\ge c\), and its outgoing edge is covered through the
interval \(J\), then
\begin{equation}
 A(T_{\rm next})
 \ge a+\mathfrak d_c(a)-\mathfrak t_J(F_c(a)).
 \label{eq:appendix-actual-local-balance}
\end{equation}
Moreover, with actual capacity slack
\[
 s(T):=1-A(T)-B(T),
\]
one has
\begin{equation}
 s(T)\ge\mathfrak d_c(A(T)).
 \label{eq:appendix-slack-release}
\end{equation}
\end{proposition}

\begin{proof}
By definition,
\[
 \widehat g_{c,J}^\vee(a)
 =\mathcal R_J(F_c(a))
 =1-e_J(F_c(a)).
\]
Insert and subtract \(F_c(a)\):
\[
 1-e_J(F_c(a))
 =1-F_c(a)-\bigl(e_J(F_c(a))-F_c(a)\bigr),
\]
which is \eqref{eq:appendix-local-balance}.  The actual statement is the
generalized handoff lemma.  Finally, the actual capped-map bound gives
\[
 B(T)\le F_c(A(T)),
\]
so
\[
 s(T)
 \ge1-A(T)-F_c(A(T))
 =\mathfrak d_c(A(T)).
\]
\end{proof}

For \(J=[L,U]\), the subsidy is explicitly
\begin{equation}
 \mathfrak t_{[L,U]}(p)=
 \begin{cases}
 0,&p<L,\\
 \max\{p,U\}-p,&p\ge L.
 \end{cases}
 \label{eq:appendix-subsidy-explicit}
\end{equation}
Thus the center is charged only for the extension connected to the incoming
trace.  Charging the full length \(U-L\) would be incorrect when \(p<L\).

\subsection{Path balance and cyclic balance}

Let \(T_m,\ldots,T_n\) be consecutive actual V triangles.  On the outgoing
edge of \(T_i\), let \(J_i\) be the external interval and put
\[
 \tau_i:=\mathfrak t_{J_i}(B_i).
\]
Coverage gives
\begin{equation}
 A_{i+1}-A_i
 \ge s(T_i)-\tau_i.
 \label{eq:appendix-signed-step}
\end{equation}
Indeed,
\[
 A_{i+1}\ge1-e_{J_i}(B_i)
 =1-B_i-\tau_i
 =A_i+s(T_i)-\tau_i.
\]
Summing yields the path balance
\begin{equation}
 A_{n+1}-A_m
 \ge\sum_{i=m}^n s(T_i)-\sum_{i=m}^n\tau_i.
 \label{eq:appendix-path-balance}
\end{equation}
On a full cycle,
\begin{equation}
 \sum_{i=0}^5s(T_i)
 \le\sum_{i=0}^5\tau_i,
 \label{eq:appendix-cycle-balance}
\end{equation}
with a strict inequality when at least one open handoff is strict after the
external intervals have been accounted for.

For a center-free chain of formal exact transfers
\[
 z_i=G_{c_i}(z_{i-1}),
\]
Definition~\ref{def:appendix-release-subsidy} gives the exact telescope
\begin{equation}
 z_r-z_0
 =\sum_{i=1}^r\mathfrak d_{c_i}(z_{i-1}).
 \label{eq:appendix-exact-telescope}
\end{equation}
This identity is the additive form of \(g\)-composition.

\begin{corollary}[Rank-zero all-Vd0 discharge]
\label{cor:appendix-rank-zero-discharge}
If six open nonsupercritical V triangles cover all six boundary edges without
an external positive-length trace on any handoff, then no such configuration
exists.
\end{corollary}

\begin{proof}
Every edge gives \(B_i+A_{i+1}>1\).  Summation gives
\[
 \sum_i(A_i+B_i)>6,
\]
contrary to \(A_i+B_i\le1\) for all \(i\).
\end{proof}

\subsection{Paired endpoint discharge}

\begin{lemma}[Endpoint release equivalence]
\label{lem:appendix-endpoint-release}
For \(a_L,a_R,c_L,c_R\in[0,1]\), put
\[
 u_L=F_{c_L}(a_L),
 \qquad
 u_R=F_{c_R}(a_R).
\]
Then
\begin{equation}
 u_L+u_R<1
 \quad\Longleftrightarrow\quad
 \mathfrak d_{c_L}(a_L)+\mathfrak d_{c_R}(a_R)
 >1-a_L-a_R.
 \label{eq:appendix-endpoint-release}
\end{equation}
The same equivalence holds with weak inequalities.
\end{lemma}

\begin{proof}
Use
\[
 u_L=1-a_L-\mathfrak d_{c_L}(a_L),
 \qquad
 u_R=1-a_R-\mathfrak d_{c_R}(a_R),
\]
and rearrange.
\end{proof}

Consequently the one-side theorem, the common CE2 two-endpoint theorem, and
the exact \(407X\) four-label theorem all have the same global form: two
endpoint diagonals release more charge than the tax left by their incoming
boundary obligations.  The strict remainder enters a center-free path and
exceeds its total nonsupercritical capacity.

\subsection{Affine deposits and threshold reservoirs}

For \(0<d<1-\sqrt3/2\), define the nonnegative threshold deposit
\begin{equation}
 \mathfrak q_d^{\rm th}(x)=
 \begin{cases}
 0,&x\le e(d),\\
 \max\{0,1-e(d)-x\},&x>e(d).
 \end{cases}
 \label{eq:appendix-threshold-deposit}
\end{equation}

\begin{lemma}[Threshold reservoir]
\label{lem:appendix-threshold-reservoir}
For all \(x\in[0,1]\),
\begin{equation}
 G_{1-d}(x)
 \ge x+\mathfrak q_d^{\rm th}(x).
 \label{eq:appendix-threshold-reservoir}
\end{equation}
If \(x>e(d)\) and \(e(d)<Q\), then
\[
 G_{1-d}(x)>1-Q.
\]
\end{lemma}

\begin{proof}
Extensivity gives \(G_{1-d}(x)\ge x\).  When \(x>e(d)\), the exact threshold
theorem gives \(G_{1-d}(x)\ge1-e(d)\).  Taking the larger of these two lower
bounds proves the first assertion.  The second follows from
\(1-e(d)>1-Q\).
\end{proof}

Suppose an affine coefficient \(\lambda\) is certified on an interval
\(I\) by
\[
 G_{1-d}(x)\ge x+\lambda(x-d)
 \qquad(x\in I).
\]
Combining this with extensivity gives the nonnegative affine deposit
\begin{equation}
 G_{1-d}(x)
 \ge x+\mathfrak q_{d,\lambda}^{\rm aff}(x),
 \qquad
 \mathfrak q_{d,\lambda}^{\rm aff}(x)
 :=\max\{0,\lambda(x-d)\}.
 \label{eq:appendix-affine-deposit}
\end{equation}
Thus a negative affine expression is never interpreted as negative released
charge; the identity wire is used in that subrange.

\begin{proposition}[CE2 one-hit discharge]
\label{prop:appendix-ce2-one-hit-discharge}
In the CE2 one-gap all-Vd0 state, one of the \(\alpha\)- or
\(\delta\)-threshold reservoirs raises the five-V-triangle chain above the
terminal target.
\end{proposition}

\begin{proof}
The signed CE2 calculation proves
\[
 X>e(\alpha),
 \qquad
 \min\{e(\alpha),e(\delta)\}<H.
\]
If \(e(\alpha)<H\), extensivity transports an input above \(e(\alpha)\) to
the \(\alpha\)-slot, and Lemma~\ref{lem:appendix-threshold-reservoir} raises
the output above \(1-H\).  Otherwise
\[
 e(\delta)<H\le e(\alpha)<X,
\]
so the same argument applies at the \(\delta\)-slot.  Subsequent extensive
maps preserve the strict inequality.
\end{proof}

\begin{proposition}[CE1 deposit cascade]
\label{prop:appendix-ce1-deposit-cascade}
In the surviving CE1 selected-\(T_+\) state, the two certified affine deposits
at deficits \(\alpha\) and \(m_3\) raise the incoming charge above
\(e(\delta)\); the final \(\delta\)-threshold reservoir then raises the
output above \(1-X\).
\end{proposition}

\begin{proof}
The universal selected-\(T_+\) concavity theorem supplies the two affine lower
bounds with coefficients \(1-4\alpha\) and \(1-5m_3\) on their certified
input intervals.  Equation~\eqref{eq:appendix-affine-deposit} converts them to
nonnegative deposits.  The exact scalar estimate in the CE1 one-gap appendix
proves that the actual second iterate is greater than \(e(\delta)\).  The
threshold reservoir and \(e(\delta)<X\) give the terminal inequality.  No
claim stronger than the certified intervals is used.
\end{proof}

\subsection{Supercritical sinks and radial conversion}

The common adjacent-rescuer theorem is already an exact discharging argument.
The rescuer and center geometry force an endpoint demand
\[
 h\ge g_c^{\rm sc},
\]
while the unique strict-supercritical V triangle has outgoing reach
\[
 b<g_c^{\rm sc}.
\]
The positive charge \(h-b\) enters a four-V-triangle center-free path.  Its
required total boundary contribution is \(4+(h-b)>4\), whereas its total
nonsupercritical capacity is at most four.

The adjacent and nonadjacent Vd terminals use a second kind of local law:
boundary charge is converted into radial separation.  The adjacent terminal
uses the quarter envelope
\[
 H\longmapsto H/4
\]
and the strict center comparison \(\delta<H/4\).  The nonadjacent terminal
uses
\[
 \min\{A,H\}>d_\tau^C
\]
and the Vd corner bound on the vertex-side radial endpoint.  These conversion
laws are type-specific and cannot be obtained from the ordinary
nonsupercritical \(g\)-map without the Vd hypotheses.

\subsection{Impossibility of a fixed universal diagonal charge}

\begin{proposition}[No positive state-independent charge]
\label{prop:appendix-no-fixed-charge}
There is no function \(q:[0,1]\to(0,\infty)\) such that
\[
 G_c(a)\ge a+q(c)
\]
for every \(a,c\in[0,1]\).  More generally, for any increasing function
\(\Phi:[0,1]\to\mathbb R\), there is no positive function \(q(c)\) such that
\[
 \Phi(G_c(a))-\Phi(a)\ge q(c)
\]
throughout the exact domain.
\end{proposition}

\begin{proof}
For every \(c\in[0,1]\), the exact capped map has a Full state with
\(a+b=1\) and \(c\le\max\{a,b\}\).  Concretely, use
\[
 (a,b)=(1-c,c)\quad(c\le1/2),
 \qquad
 (a,b)=(c,1-c)\quad(c\ge1/2).
\]
At that state \(F_c(a)=1-a\), hence \(G_c(a)=a\) and
\(\mathfrak d_c(a)=0\).  Both asserted positive lower bounds are therefore
impossible.
\end{proof}

The same example rules out a fixed positive split of diagonal charge between
the two incident boundary edges.  The endpoint theorems can force only a
joint release after both incoming endpoint states have been specified.

\subsection{Why gap length and boundary balance alone are insufficient}

A singleton V-gap is nonempty but has length zero.  It is active because two
open incident traces do not cover their common endpoint.  Therefore a charge
system recording only uncovered length cannot distinguish a singleton active
gap from a genuinely gap-free handoff.  The connected residual operator and
the active-gap convention are indispensable.

Boundary balance alone also does not eliminate the \(N_+=1\), rank-zero
all-Vd0 state.  To see the algebraic obstruction, put
\[
 \varepsilon=\frac1{200},
\]
\[
 (x_0,x_1,x_2,x_3,x_4,x_5)
 =\left(\frac{55}{100},\frac{53}{100},\frac{51}{100},
        \frac{49}{100},\frac{47}{100},\frac{45}{100}\right),
\]
and define cyclically
\[
 B_i=x_i,
 \qquad
 A_i=1-x_{i-1}+\varepsilon.
\]
Then every boundary handoff is strict:
\[
 B_i+A_{i+1}=1+\varepsilon>1.
\]
Exactly one role is supercritical:
\[
 A_0+B_0=1+\varepsilon+\frac{10}{100}>1,
\]
whereas
\[
 A_i+B_i=1+\varepsilon-\frac2{100}<1
 \qquad(1\le i\le5).
\]
The exceptional pair also satisfies the elementary endpoint-distance bound
\[
 A_0^2+A_0B_0+B_0^2<1.
\]
Thus skeleton handoffs, six capacity signs, and the basic boundary diameter
constraint admit consistent numerical data.  The active proof must use the
six radial witnesses and three asymmetric frontier witnesses of the
nine-point obstruction.

\subsection{Audit of successful and unsuccessful variants}

\begin{table}[H]
\centering
\scriptsize
\renewcommand{\arraystretch}{1.12}
\begin{tabular}{@{}p{.30\textwidth}p{.16\textwidth}p{.43\textwidth}@{}}
\toprule
variant&status&reason\\
\midrule
fixed charge \(q(c)>0\)&fails&Full states have zero release\\
fixed left/right split&fails&either endpoint may have zero release; only the
paired release is controlled\\
constant increment after a potential change&fails&every increasing potential
has zero increment at a Full state\\
static order-free charge sum&fails&threshold activation depends on the arriving
state\\
state-dependent release \(\mathfrak d_c(a)\)&works&it is exactly the capped
\(g\)-transfer in additive form\\
release minus connected center subsidy&works&Proposition
\ref{prop:appendix-local-balance} is an identity\\
paired endpoint discharge&works&equivalent to the exact endpoint sum
inequalities\\
CE2 one-hit threshold&works&one of two reservoirs must fire\\
CE1 affine-deposit cascade&works&two certified deposits cross the final
threshold\\
adjacent-rescuer surplus&works&the source--sink difference exceeds the
four-role path capacity\\
pure boundary discharge for rank-zero \(N_+=1\)&fails&the displayed formal
boundary data are feasible; interior witnesses are needed\\
scalar replacement of the two-chart Vd1 pair&fails&orientation, adjacent-arm
exclusion, and radial preservation are not scalar data\\
\bottomrule
\end{tabular}
\caption{Audit of discharging variants for Strategy~2.}
\label{tab:appendix-discharging-audit}
\end{table}
